% =============================================================================
\chapter{Estrutura do repositório}
\label{ap:repositorio}
% =============================================================================

A árvore abaixo mostra o que importa em cada pasta. A separação entre o que é
produto (o que roda quando alguém usa o Morphe) e o que não é está também no
\arq{README.md} da raiz.

\medskip
% Arvore do repositorio (estado de 24/09/2026, com o ADC). Pacote dirtree.
\renewcommand*\DTstylecomment{\small\rmfamily\itshape\color{black!65}}
\renewcommand*\DTstyle{\small\ttfamily}
\dirtree{%
.1 morphe/.
.2 README.md\DTcomment{ponto de entrada}.
.2 morphe-up.sh\DTcomment{prepara a placa: FPGA, servidor, marca}.
.2 instala-quartus.sh\DTcomment{Quartus 20.1 para todas as contas de um PC}.
.2 instala-morphe.sh\DTcomment{o aplicativo em /opt/morphe, com atalho}.
.2 placas.conf\DTcomment{as placas do laboratório}.
.2 deploy.sh, gera\_proveniencia.sh, PROVENIENCIA.sha256.
.2 Python/\DTcomment{o aplicativo}.
.3 morphe\_app.py\DTcomment{janela principal}.
.3 *\_window.py\DTcomment{uma janela por operação}.
.3 morphe\_protocol.py\DTcomment{cliente TCP}.
.3 dsp\_core.py, blocos.py\DTcomment{ponto fixo, referência, sinais longos}.
.3 fir\_design.py, iir\_design.py\DTcomment{projetistas de filtros}.
.3 aquisicao.py, sinal\_arquivo.py\DTcomment{ADC e sinais de arquivo}.
.3 ferramentas/\DTcomment{testes contra a placa, figuras, vetores}.
.2 C/\DTcomment{o servidor da placa}.
.3 morphe\_server.c, morphe\_protocol.h, morphe\_config.h.
.3 hps\_0.h\DTcomment{gerado do Platform Designer}.
.3 autostart/.
.2 Quartus/\DTcomment{o hardware}.
.3 ghrd\_top.v\DTcomment{topo: instancia os aceleradores}.
.3 conv1d.v, fft\_wrapper.v, iir\_*.v\DTcomment{os aceleradores}.
.3 adc\_captura.v\DTcomment{o controlador do ADC}.
.3 soc\_system.qsys, soc\_system.qsf.
.3 tb\_*.v, roda\_tb\_*.sh\DTcomment{testbenches (IIR e ADC)}.
.3 gen\_hps\_header.py\DTcomment{gera o hps\_0.h}.
.3 output\_files/\DTcomment{o bitstream}.
.2 exemplos/\DTcomment{pacotes .mrph e sinais de exemplo}.
.2 docs/\DTcomment{documentação técnica, diário e este manual}.
.2 legado/\DTcomment{do projeto original, fora do produto}.
}

\medskip

Três regras sobre esta estrutura:

\begin{itemize}
  \item \textbf{os \arq{iir_*.v} não saem de \arq{Quartus/}}: não estão
    listados no \arq{soc_system.qsf}, e o Quartus os acha por estarem na pasta
    do projeto;
  \item \textbf{o \arq{C/hps_0.h} e o \arq{Quartus/soc_system.sopcinfo} são
    gerados}, e vão juntos com o bitstream que os gerou;
  \item \textbf{as ferramentas rodam de dentro de \arq{Python/}}, porque o
    \arq{ferramentas/_caminho.py} acrescenta a pasta de cima ao caminho de
    importação; por exemplo:\\
    \texttt{python3 ferramentas/testa\_iir\_hw.py <ip>}.
\end{itemize}

\begin{table}[htbp]
  \centering
  \small
  \caption{A documentação técnica, em \arq{docs/} e na raiz.}
  \label{tab:docs}
  \begin{tabularx}{\textwidth}{@{}lX@{}}
    \toprule
    \textbf{Arquivo} & \textbf{Assunto} \\
    \midrule
    \arq{PREPARACAO.md} & preparar a placa, instalar para a turma, Quartus nos
      computadores \\
    \arq{docs/GERENCIAMENTO-PLACAS.md} & várias placas e vários clientes \\
    \arq{docs/PRECISAO-NUMERICA.md} & os formatos de ponto fixo e os erros
      medidos \\
    \arq{docs/ADC.md} & o conversor da placa e como a captura funciona \\
    \arq{docs/COMPILAR-IIR.md}, \arq{docs/COMPILAR-ADC.md} & recompilar o
      bitstream, passo a passo \\
    \arq{docs/NOVA-PLACA.md} & pôr uma placa nova em serviço \\
    \arq{docs/RESSALVAS.md} & o que foi corrigido no projeto original, e por
      quê \\
    \arq{docs/LINHA-DE-BASE-PASSOS.md} & a medição de usabilidade antes das
      mudanças \\
    \arq{docs/DIARIO.md} & o diário do estágio: o que foi feito, medido e
      validado, dia a dia \\
    \arq{docs/manual/} & este manual \\
    \bottomrule
  \end{tabularx}
\end{table}

% =============================================================================
\chapter{Referência rápida}
\label{ap:referencia}
% =============================================================================

Os comandos do dia a dia. Os da estação rodam de \arq{/opt/morphe}, na conta de
quem vai preparar a aula; nunca com \texttt{sudo}.

\newcommand{\cmd}[1]{\par\smallskip\noindent\hspace*{1em}{\ttfamily\small #1}\par\smallskip}

\section*{Preparar as placas (estação)}

\cmd{cd /opt/morphe \&\& ./morphe-up.sh -{}-todas}

Só uma placa:

\cmd{./morphe-up.sh -{}-cable 'DE-SoC [1-2]'}
\cmd{./morphe-up.sh -{}-cable 'DE-SoC [1-3]'}

Ver o que está no ar, e se os \emph{tethers} estão vivos:

\cmd{./morphe-up.sh -{}-status}
\cmd{pgrep -af quartus\_pgm}

\section*{Abrir o aplicativo (qualquer computador)}

\cmd{cd /opt/morphe/Python \&\& python3 morphe\_app.py}

\section*{Uma vez por placa e por conta}

\cmd{./morphe-up.sh -{}-setup-ssh -{}-board <ip>}

\section*{Atualizar (estação e qualquer computador)}

Numa conta com \texttt{sudo}, ou na conta dona do \arq{/opt/morphe}:

\cmd{/opt/morphe/atualizar.sh}

\noindent Na estação, se o servidor ou o bitstream mudaram,
\arq{./morphe-up.sh} \opt{todas} de novo. Computador instalado antes de
01/10/2026, uma única vez:

\cmd{sudo -u \$(stat -c \%U /opt/morphe) git -C /opt/morphe pull -{}-ff-only}

\section*{Os outros computadores}

Instalar, numa conta com \texttt{sudo} (uma vez por computador):

\cmd{git clone -b estagio/v1.1-1024pontos https://github.com/nicassiosantos/morphe.git /tmp/morphe-inst}
\cmd{sudo /tmp/morphe-inst/instala-morphe.sh}

Atualizar: \arq{/opt/morphe/atualizar.sh}, como acima.

\noindent Só conferir, sem mudar nada: \texttt{sudo /opt/morphe/instala-morphe.sh -{}-verificar}.

\section*{O que cada placa está fazendo}

\cmd{python3 /opt/morphe/Python/ferramentas/estado\_placas.py -{}-seguir}

\section*{Conferir um clone}

\cmd{sha256sum -c PROVENIENCIA.sha256}
\cmd{python3 Quartus/gen\_hps\_header.py -{}-check}

\section*{Testar contra a placa (de dentro de \texttt{Python/})}

\cmd{python3 ../Quartus/morphe\_ping.py <ip>}
\cmd{python3 ferramentas/testa\_ifft\_roundtrip.py <ip>}
\cmd{python3 ferramentas/testa\_iir\_hw.py <ip>}
\cmd{python3 ferramentas/testa\_blocos.py -{}-placa <ip>}
\cmd{python3 ferramentas/testa\_adc.py <ip> basico}
\cmd{python3 ferramentas/testa\_concorrencia.py -{}-modo mesma -{}-clientes 6}

\section*{Encerrar}

\cmd{./morphe-up.sh -{}-down}

\noindent encerra os \emph{tethers}; em uma hora as placas param. Com
\opt{stop-server}, para também os servidores.

\section*{Quando algo dá errado}

\begin{center}
\small
\begin{tabular}{@{}ll@{}}
  \toprule
  \textbf{Sintoma} & \textbf{Primeiro passo} \\
  \midrule
  tempo esgotado, FFT errada & \texttt{pgrep -af quartus\_pgm}; \arq{./morphe-up.sh} \\
  ``FPGA não preparada'' & \arq{./morphe-up.sh} \\
  aplicativo não acha placa & \arq{estado_placas.py};
    \arq{./morphe-up.sh} \opt{todas} \\
  resultado inexplicável & \texttt{gen\_hps\_header.py -{}-check} \\
  aplicativo não abre & remover NumPy e Matplotlib do \texttt{pip -{}-user} \\
  \bottomrule
\end{tabular}
\end{center}

\noindent As tabelas completas estão nos capítulos~\ref{cap:placas}
e~\ref{cap:laboratorio}.
